\section{問題提起}

既存のオブジェクト検出アルゴリズム（YOLOシリーズやSSD等）は人物個体の検出には長けているものの、人物同士の相互関係を捉えてグループ化する機能は直接備えていない。そのため、待機列映像からグループを自動検出するシステムを構築・実用化するにあたっては、単純な人物検出手法の適用だけでは解決できない以下の4つの技術的・環境的課題が存在する。

\begin{enumerate}
    \item \textbf{密集と遮蔽（オクルージョン）による検出漏れおよび重複}：
    待機列内では人物同士の距離が極めて近く、人物の一部や全身が他者によって遮られる現象が頻発する。これにより、標準的なバウンディングボックスによるオブジェクト検出では、一人の人物に対して複数の検出枠が乱立する、あるいは密集部位の人物が完全に無視されるといった問題が生じる。

    \item \textbf{複雑な待機列構造におけるグループ境界の曖昧性}：
    グループの形成パターンは物理的な近接度のみに依存しない。待機列の折返しや壁沿いの待機、スマートフォンの操作に伴う体の向きの変化、前後列に跨る会話など、単一フレームの二次元距離情報のみでは「単なる前後の他人」と「会話中の同一グループ」の判別が困難である。

    \item \textbf{一時的な通行人やノイズによる誤検出}：
    待機列の周辺を通行する人物や、画角の端を通過する一時的な対象が待機列構成員として認識されることで、グループ数および待ち時間算出の精度が著しく低下する。

    \item \textbf{Webエッジ環境における実行リソースとOOM（Out of Memory）の制約}：
    専用サーバーを介さずブラウザ上で高頻度な推論処理（オブジェクト検出・ポーズ推定・複数フレーム追跡など）を実行する場合、メモリの無制限な蓄積や高頻度な推論によるブラウザクラッシュ（OOM）の危険性が高まる。また、UIスレッドのブロックを避けつつリアルタイム性を維持するための並列処理・構造単純化が不可欠となる。
\end{enumerate}

これらの複合的な課題を同時に克服するアルゴリズムの設計と、軽量かつ安定したシステムアーキテクチャの確立が不可欠となっている。
